\chapter{Certificats de dépôt et double cotation}\label{ch:hf:depositary-receipts-and-dual-listings}

Les actions d'une société se négocient à Londres et ses certificats de dépôt à New York. Pendant deux heures, les deux marchés sont ouverts et l'un n'est qu'une conversion monétaire de l'autre~; pour le
reste de la journée new-yorkaise, l'un d'eux est évalué par un teneur de marché qui ne peut qu'estimer. Dans la simulation de ce chapitre, la fair value du certificat est
connue à 2 points de base tant que le marché d'origine est ouvert, à 74 une fois celui-ci fermé, et à 94 sans proxy pour ce marché~; une conversion par le
dépositaire ne rapporte que lorsque l'écart dépasse 29 points de base.

\section{Certificats, ratios et dépositaire}

Un certificat de dépôt (livre~1, chapitre~17) est émis par une banque dépositaire contre des actions qu'elle détient en conservation sur le marché d'origine. Chaque certificat
représente un nombre fixe d'actions.

\begin{definition}[Ratio du certificat de dépôt]\label{def:hf:depositary-receipts-and-dual-listings:ratio}
Le \emph{ratio du certificat de dépôt}\index{ratio du certificat de dépôt} est le nombre d'actions du marché d'origine que représente un certificat (ou son inverse,
selon ce qu'indique le dépositaire)~; avec lui et le taux de change, la valeur de parité d'un certificat vaut le prix du marché d'origine multiplié par le ratio et par le prix de la monnaie dans
la monnaie du certificat.
\end{definition}

\begin{definition}[Émission de certificats de dépôt]\label{def:hf:depositary-receipts-and-dual-listings:issuance}
L'\emph{émission de certificats de dépôt}\index{emission@émission de certificats de dépôt} est la création de nouveaux certificats par le dépositaire contre des actions remises
à sa garde~; l'annulation en est l'inverse, des certificats restitués contre les actions qu'ils représentent. Les deux supportent les frais du dépositaire par certificat et
se règlent selon les cycles de règlement-livraison des deux marchés.
\end{definition}

L'émission et l'annulation jouent pour les certificats le rôle que la création et le rachat jouent pour les fonds indiciels cotés (ETF) (chapitre~12)~: elles permettent à un arbitragiste de
déplacer leur inventaire entre les deux marchés lorsque les prix s'écartent. Une double cotation, soit les mêmes actions cotées sur deux bourses, n'exige pas de dépositaire~: les
actions passent d'un registre à l'autre, avec le même effet économique et ses propres frictions.

\begin{definition}[Pré-émission]\label{def:hf:depositary-receipts-and-dual-listings:prerelease}
La \emph{pré-émission}\index{pré-émission} (pre-release) est l'\omterm{def:hf:depositary-receipts-and-dual-listings:issuance}{émission de certificats de dépôt} avant le dépôt des actions correspondantes, autorisée par les
conventions du dépositaire seulement si le courtier qui les reçoit a conclu un accord avec le dépositaire et si le courtier ou son client détient les
actions correspondantes.
\end{definition}

\begin{dated}{2026-09}{Pré-émission sanctionnée, et frais d'un dépositaire}\label{dat:hf:depositary-receipts-and-dual-listings:prerelease}
Le 26 décembre 2018, la SEC annonça que JPMorgan paierait plus de 135~millions de dollars pour régler des accusations de traitement irrégulier de
certificats de dépôt américains pré-émis~: elle avait remis des certificats à des courtiers dans des milliers de transactions de \omterm{def:hf:depositary-receipts-and-dual-listings:prerelease}{pré-émission} alors que ni le courtier ni ses clients
ne détenaient les actions étrangères nécessaires à leur couverture. La SEC indiqua qu'elle avait alors tenu pour responsables les quatre banques dépositaires~; ses communiqués de 2018 annoncèrent
des règlements d'environ 75~millions de dollars avec des filiales de Deutsche Bank, plus de 54~millions de dollars avec BNY Mellon et plus de 38~millions de dollars avec Citibank. Un
barème des frais d'un dépositaire, tel que NIO le communiqua dans son rapport annuel 2018, prélevait jusqu'à 0.05~\$ par certificat émis et jusqu'à 0.05~\$ par
certificat annulé.
\end{dated}

\section{Arbitrage pendant les heures de chevauchement}

Tant que les deux marchés sont ouverts, le certificat est un portefeuille de trois prix~: les actions d'origine, la monnaie et le certificat lui-même. Gagnon et Karolyi
ont comparé les cours intrajournaliers des certificats aux cours synchrones du marché d'origine, ajustés de la monnaie, pour 506 valeurs cotées sur deux places issues de 35 pays~: les écarts
à la parité s'élevaient en moyenne à 4.9 points de base, un niveau économiquement faible, mais ils étaient volatils et pouvaient être importants, et ils croissaient avec les indicateurs des coûts de détention qui
entravent l'arbitrage. Un teneur de marché des certificats, pendant le chevauchement, cote le certificat à partir du cours en temps réel du marché d'origine (le \omterm{def:hf:fair-value:cross}{prix juste inter-instruments} du chapitre~2)
et se couvre avec les actions d'origine et la monnaie~; un arbitragiste prend position sur le certificat lorsque son écart à la parité dépasse l'aller-retour sur les trois instruments.

Dans la simulation du chapitre, avec 2 points de base de bruit dans le mid d'origine et 4 dans celui du certificat, l'écart entre le certificat et la parité a un
écart type de 4.5 points de base et dépasse un aller-retour de 10 points de base dans 2.7~\% des minutes de chevauchement~: un arbitrage qui existe, rarement,
et revient à celui qui le voit en premier.

\section{Valorisation hors des heures de chevauchement}

Lorsque le marché d'origine ferme, le certificat continue de se négocier, et son teneur de marché a besoin d'une fair value sans le cours d'origine. Le dernier cours d'origine, corrigé de la
monnaie, est la réponse naïve~; le dernier cours d'origine corrigé par un proxy encore négocié (un contrat à terme sur l'indice d'origine, un fonds
coté sur le marché d'origine, à côté du certificat) et par la monnaie est préférable~; les négociations du certificat lui-même et les cours d'autres certificats du même pays,
apportent une information que le proxy n'a pas.

\begin{omfigure}
\begin{tikzpicture}
\begin{axis}[width=11cm,height=5.2cm,xlabel={\small minutes depuis le début de la séance du certificat},ylabel={\small erreur de cotation (pb)},xmin=0,xmax=390,
  ymin=0,ymax=145,tick label style={font=\footnotesize},grid=major,grid style={gray!25},table/col sep=comma,
  legend style={font=\scriptsize,at={(0.5,1.03)},anchor=south}]
\fill[gray!12] (axis cs:0,0) rectangle (axis cs:120,145);
\addplot[omDef,thick,dashed,mark=square*,mark size=1.1pt] table[x=minute,y=home]{figdata/hft/14-depositary-receipts-and-dual-listings/error.csv};
\addplot[omThm,thick,mark=*,mark size=1.1pt] table[x=minute,y=proxy]{figdata/hft/14-depositary-receipts-and-dual-listings/error.csv};
\legend{{fermeture du marché d'origine et monnaie},{fermeture du marché d'origine, monnaie et proxy}}
\node[font=\scriptsize,anchor=south,align=center] at (axis cs:60,20) {les deux marchés\\ouverts};
\end{axis}
\end{tikzpicture}

\omcaption{Erreur quadratique moyenne de la fair value du teneur de marché des certificats par rapport à la valeur vraie du certificat, par quart d'heure de la séance du certificat~;
le marché d'origine est ouvert pendant les deux premières heures (en grisé). Après sa fermeture, la fair value est le dernier mid d'origine corrigé de la monnaie, avec ou
sans le mouvement d'un proxy du marché d'origine (volatilité quotidienne de 1~\% pour le marché, de 1.2~\% pour la part propre de la société, de 0.2~\% pour la
base du proxy). Trois paires simulées de 40 séances. Données~: \texttt{hf\_dr.by\_minute}.}\label{fig:hf:depositary-receipts-and-dual-listings:error}
\end{omfigure}

L'erreur croît avec la racine carrée du temps écoulé depuis la fermeture du marché d'origine (la \cref{fig:hf:depositary-receipts-and-dual-listings:error})~: le proxy
retire la part du marché et laisse celle propre à la société, que nul proxy ne peut porter. Un teneur de marché qui cote le certificat à ces heures élargit sa
fourchette avec le temps écoulé depuis la fermeture du marché d'origine, biaise ses cotations selon sa prévision et couvre la part du marché dans le proxy.

Froot et Dabora ont étudié un cas plus tranché, celui de sociétés qui mettent leurs flux de trésorerie en commun mais cotent des actions jumelles sur des marchés différents~: le prix relatif des jumelles
évoluait avec les marchés où chacune se négociait le plus. Le prix relatif d'une jumelle montait lorsque le marché sur lequel elle se négociait relativement davantage montait~; aucune explication testée
(dividendes, dépenses, droits de vote, monnaie, date de détachement, impôts) n'expliquait pleinement cette évolution.

\section{Conversion~: émission, annulation et ses frictions}

La conversion déplace des actions d'un marché à l'autre à un coût~: les frais du dépositaire par certificat, les fourchettes des jambes, la monnaie et le financement sur le
délai de règlement-livraison. Pour un certificat de 25~\$, des frais de 0.05~\$ représentent 20 points de base~; avec 6 points de base de coûts de négociation et deux jours de financement à 5~\%, l'écart doit
dépasser 28.8 points de base avant que la conversion ne rapporte (le \cref{lst:hf:depositary-receipts-and-dual-listings:conversion}). Un teneur de marché convertit donc
pour déplacer un inventaire qu'il ne pourrait pas déboucler par la négociation, et non pour capter l'écart courant. Sur les marchés soumis à des plafonds de détention étrangère, la conversion en
certificats peut être plafonnée, et le certificat se négocie alors avec une prime durable que nul arbitrage ne peut résorber.

\section{Fiches de stratégie}

\begin{strategyfile}{Certificat contre actions ordinaires pendant les heures de chevauchement}\label{strat:hf:depositary-receipts-and-dual-listings:overlap}
\sfield{Qui paie, et pourquoi} Les investisseurs dans le certificat ou dans les actions d'origine qui négocient sans surveiller l'autre cotation, et les fournisseurs de liquidité
dont les cotations suivent l'autre marché avec retard.
\sfield{Instruments et plateformes} Le certificat sur une bourse américaine, les actions ordinaires d'origine, la monnaie.
\sfield{Signal} Le cours du certificat comparé au cours d'origine multiplié par le ratio et par la monnaie, au-delà du coût d'aller-retour des trois jambes.
\sfield{Taille et exécution} Vendre le côté le plus cher et acheter le moins cher~; ou coter le certificat à partir du marché d'origine et couvrir les exécutions sur place et en
monnaie~; conserver la position jusqu'à ce que l'écart se referme ou convertir lorsque l'inventaire s'accroît.
\sfield{Coûts} Trois fourchettes, frais sur les deux bourses, droits de timbre au marché d'origine lorsqu'ils s'appliquent, fourchette de la monnaie.
\sfield{Comment elle disparaît} Les écarts sont en moyenne de quelques points de base (4.9 dans l'échantillon de Gagnon et Karolyi) et vont au lien le plus rapide entre les deux marchés.
\sfield{Horizon, capacité, infrastructure} De la milliseconde à la minute pendant le chevauchement~; liens entre les centres de données des deux marchés.
\sfield{Backtest rigoureux} Prix synchrones sur une même horloge~; prix exécutable de la monnaie~; règlement-livraison et droit de timbre du marché d'origine.
\sfield{Sources} Gagnon et Karolyi (2010).
\end{strategyfile}

\begin{strategyfile}{Fixer le prix du certificat lorsque le marché d'origine est fermé}\label{strat:hf:depositary-receipts-and-dual-listings:closed}
\sfield{Qui paie, et pourquoi} Des investisseurs qui négocient la société pendant leurs propres heures et paient pour l'immédiateté lorsque le marché d'origine ne peut pas la fournir.
\sfield{Instruments et plateformes} Le certificat~; un proxy du marché d'origine (contrat à terme sur indice, fonds pays)~; la monnaie.
\sfield{Signal} La clôture du marché d'origine déplacée par le proxy et par la monnaie, mise à jour par les transactions sur le certificat lui-même et par les autres certificats du même marché.
\sfield{Taille et exécution} Coter avec une fourchette qui s'élargit avec le temps écoulé depuis la clôture du marché d'origine~; couvrir la part de marché dans le proxy~; conserver la
part propre à la société jusqu'à l'ouverture du marché d'origine.
\sfield{Coûts} La fourchette et la base du proxy~; le risque spécifique à la société pendant la nuit.
\sfield{Comment elle disparaît} Les annonces de la société pendant que le marché d'origine est fermé~: dans la simulation, l'erreur avec le proxy atteint encore 109 points de base à la
fin de la séance.
\sfield{Horizon, capacité, infrastructure} Des heures~; capacité tirée du volume du certificat et de l'appétit du teneur de marché pour le risque de nuit.
\sfield{Backtest rigoureux} Le prix du proxy à chaque cotation~; les annonces de la société horodatées au moment de leur publication~; l'ouverture du marché d'origine, et non la clôture, comme sortie.
\sfield{Sources} Ce chapitre~; Froot et Dabora (1999) sur la façon dont le lieu des transactions influe sur les prix.
\end{strategyfile}

\begin{strategyfile}{Arbitrage de conversion par le dépositaire}\label{strat:hf:depositary-receipts-and-dual-listings:conversion}
\sfield{Qui paie, et pourquoi} Une demande concentrée sur une seule cotation (inclusion dans un indice, flux d'un fonds local) qui pousse celle-ci à s'écarter de l'autre.
\sfield{Instruments et plateformes} Les certificats, les actions ordinaires, l'émission et l'annulation de certificats par le dépositaire, la monnaie.
\sfield{Signal} Un écart au-delà du seuil de conversion~: les frais par certificat en proportion de son cours, les coûts des jambes et le financement pendant le
délai de règlement-livraison.
\sfield{Taille et exécution} Acheter la cotation bon marché, convertir, livrer dans la cotation chère~; dimensionner la position selon la persistance de l'écart, et non selon son pic.
\sfield{Coûts} Les frais du dépositaire (jusqu'à 0.05~\$ par certificat dans chaque sens selon le barème publié), la conservation, les taxes sur les transferts, le financement.
\sfield{Comment elle disparaît} Les plafonds de détention qui limitent la conversion~; et la tentation de la \omterm{def:hf:depositary-receipts-and-dual-listings:prerelease}{pré-émission}, dont les affaires de la SEC de 2018 montrent qu'elle franchit une limite légale.
\sfield{Horizon, capacité, infrastructure} Des jours~; capacité tirée du traitement du dépositaire et de la taille de l'écart.
\sfield{Backtest rigoureux} Dates de règlement-livraison sur les deux marchés, barème des frais en vigueur et temps réel d'achèvement de la conversion.
\sfield{Sources} Communiqués de presse de la SEC sur la \omterm{def:hf:depositary-receipts-and-dual-listings:prerelease}{pré-émission} (2018)~; rapport annuel de NIO (barème des frais).
\end{strategyfile}

\section{Tutoriel~: deux prix, une même société}\label{tut:hf:depositary-receipts-and-dual-listings:tut}

\begin{tutorial}
\textbf{Objectif.} Mesurer à quel point un teneur de marché de certificats connaît leur fair value pendant les heures d'ouverture et de fermeture, et à partir de quand la conversion devient rentable. \textbf{État final~:}
la \cref{fig:hf:depositary-receipts-and-dual-listings:error} et les chiffres du texte.
\begin{tutsteps}
\item \textbf{La paire.} \texttt{firm.drarb.Pair}~: 40 séances de 390 minutes pour le certificat, le marché d'origine ouvert pendant les 120 premières~; un certificat de deux
actions à 1.25~dollars par unité de monnaie d'origine~; la valeur de la société se décompose en une part de marché (1~\% par jour) et une part propre (1.2~\%)~; la monnaie 0.6~\%~; un proxy
suivant la part de marché, avec une base de 0.2~\% par jour.
\item \textbf{Fair values} à partir du dernier mid du marché d'origine, de la monnaie et, après la clôture du marché d'origine, du mouvement du proxy.
\omcode{code/firm/drarb/firm_drarb.py}{83}{87}{La fair value du certificat~: dernier mid du marché d'origine multiplié par le ratio et par la monnaie, déplacé par le proxy depuis la clôture du marché d'origine.}\label{lst:hf:depositary-receipts-and-dual-listings:quote}
\item \textbf{Erreurs} par rapport à la valeur vraie, par minute de la séance (\texttt{hf\_dr.errors}, \texttt{hf\_dr.by\_minute}).
\item \textbf{Conversion.} Le seuil pour un certificat de 25~\$ avec des frais de 0.05~\$.
\omcode{code/firm/drarb/firm_drarb.py}{40}{52}{Le seuil de conversion et la décision~: frais, coûts de négociation, financement sur le délai de règlement-livraison.}\label{lst:hf:depositary-receipts-and-dual-listings:conversion}
\end{tutsteps}
\textbf{Pour aller plus loin.} Mettre à jour la fair value à partir des transactions propres au certificat après la clôture du marché d'origine (filtre du chapitre~2)~; ajouter d'autres certificats
du même pays comme second proxy~; modéliser l'écart de nuit jusqu'à l'ouverture du marché d'origine.
\end{tutorial}

Les chiffres~: 2.0 points de base d'erreur quadratique moyenne tant que le marché d'origine est ouvert (le bruit propre au mid du marché d'origine), 74 après sa clôture avec le proxy
et 94 sans, qui montent à 109 et 137 à la fin de la séance. Le proxy élimine la part de marché de l'erreur et laisse celle de la société.

\section{Construction~: le module d'arbitrage sur certificats}\label{bld:hf:depositary-receipts-and-dual-listings:drarb}

\begin{build}
\bfield{Rôle} Valoriser un certificat pendant les heures d'ouverture et de fermeture, et décider des conversions.
\bfield{Interface} \texttt{dr\_fair(home, ratio, fx)}, \texttt{closed\_fair(home\_close, ratio, fx, proxy\_move, beta)},
\texttt{threshold\_bp(dr\_price, fee\_per\_dr, cost\_bp, days, rate)},
\texttt{conversion\_edge(dr, home, ratio, fx, fee\_per\_dr, cost\_bp, days, rate)}, \texttt{Pair(days, seed)} avec \texttt{quote(method)} et \texttt{error\_bp(method)}.
\bfield{Règles} La monnaie s'exprime en dollars par unité de monnaie d'origine~; les coûts de conversion incluent le financement sur le délai de règlement-livraison.
\bfield{Tests d'acceptation} \texttt{code/firm/drarb/tests/}~: fair values calculées à la main~; seuil et edges pour l'émission, l'annulation et l'absence d'opération~; la référence du
marché d'origine est la minute en cours pendant l'ouverture et la clôture ensuite~; les deux méthodes concordent pendant l'ouverture~; le proxy réduit l'erreur à marché fermé, qui
reste plus de dix fois supérieure à l'erreur à marché ouvert.
\bfield{Pour aller plus loin} Les transactions propres au certificat comme source d'information~; les écarts de nuit~; les plafonds de détention appliqués à la conversion.
\end{build}

\begin{omsources}
\item L.~Gagnon, G.~A.~Karolyi, Multi-market trading and arbitrage, \emph{Journal of Financial Economics} 97(1), 2010, 53--80.
\item K.~A.~Froot, E.~M.~Dabora, How are stock prices affected by the location of trade?, \emph{Journal of Financial Economics} 53(2), 1999,
189--216.
\item US Securities and Exchange Commission, communiqués de presse 2018-138, 2018-255, 2018-285 et 2018-306 sur les ADR pré-émis.
\item NIO Inc., rapport annuel sur le Form 20-F pour 2018, frais et commissions pour les détenteurs d'ADS.
\end{omsources}

\section{Exercices}

\begin{exercise}[$\star$]\label{exo:hf:depositary-receipts-and-dual-listings:1}
Sachant que les actions d'origine se négocient à 10 en monnaie d'origine, que le ratio est de deux actions par certificat et que la monnaie vaut 1.25~dollars, calculer la parité, puis, après la
clôture du marché d'origine, sachant qu'un proxy monte de 1~\% et que le bêta de l'action par rapport à lui vaut 1.2, calculer la fair value du certificat.
\end{exercise}

\begin{exercise}[$\star$]\label{exo:hf:depositary-receipts-and-dual-listings:2}
Calculer le seuil de conversion pour un certificat de 25~\$ avec des frais de 0.05~\$, 6 points de base de coûts de négociation et deux jours de financement à 5~\%.
\end{exercise}

\begin{exercise}[$\star$]\label{exo:hf:depositary-receipts-and-dual-listings:3}
Le certificat se négocie à 25.20~\$ et la parité est de 25.00~\$. Indiquer s'il faut émettre, annuler ou ne rien faire, et calculer l'edge par certificat.
\end{exercise}

\begin{exercise}[$\star\star$]\label{exo:hf:depositary-receipts-and-dual-listings:4}
Expliquer pourquoi l'erreur à marché fermé croît avec la racine carrée du temps, et ce qui subsiste après le proxy.
\end{exercise}

\begin{exercise}[$\star\star$]\label{exo:hf:depositary-receipts-and-dual-listings:5}
Sachant que Gagnon et Karolyi ont observé des écarts de parité de 4.9 points de base en moyenne, expliquer pourquoi un teneur de marché peut encore gagner sur les certificats.
\end{exercise}

\begin{exercise}[$\star\star$]\label{exo:hf:depositary-receipts-and-dual-listings:6}
Expliquer ce qui rendait la \omterm{def:hf:depositary-receipts-and-dual-listings:prerelease}{pré-émission} irrégulière dans les affaires de la SEC de 2018.
\end{exercise}

\begin{exercise}[$\star\star\star$]\label{exo:hf:depositary-receipts-and-dual-listings:7}
\emph{Programmation.} Relancer \texttt{hf\_dr.errors} avec la volatilité quotidienne propre à la société à 0.6~\% au lieu de 1.2~\%. Calculer la baisse de
l'erreur à marché fermé avec le proxy.
\end{exercise}

\begin{exercise}[$\star\star\star$]\label{exo:hf:depositary-receipts-and-dual-listings:8}
\emph{Trouver la faille.} «~Le certificat se négocie 40 points de base sous la parité à 21:00 heure de Londres~: il suffit d'acheter le certificat et de vendre les actions à l'ouverture de demain.~»
\end{exercise}

\section{Problème~: deux prix, une même société}

\begin{problem}[{Problème du week-end --- deux prix, une même société}]\label{pb:hf:depositary-receipts-and-dual-listings:1}
Un teneur de marché cote le certificat d'une société à New York pendant les heures où son marché d'origine est ouvert et celles où il est fermé.

\medskip\noindent\textbf{Partie I --- L'instrument.}
\begin{enumerate}
\item Définir le \omterm{def:hf:depositary-receipts-and-dual-listings:ratio}{ratio du certificat de dépôt} et la parité.
\item Définir l'\omterm{def:hf:depositary-receipts-and-dual-listings:issuance}{émission de certificats de dépôt} et leur annulation.
\item Définir la \omterm{def:hf:depositary-receipts-and-dual-listings:prerelease}{pré-émission} (pre-release) et la condition qui la rend valable.
\item Résumer l'encadré daté.
\end{enumerate}

\medskip\noindent\textbf{Partie II --- Prix.}
\begin{enumerate}[resume]
\item Que mesurent Gagnon et Karolyi, et que trouvent-ils~?
\item Que trouvent Froot et Dabora au sujet des actions jumelles~?
\item Écrire le fair value pendant les heures de fermeture du marché d'origine.
\item Quelle information la négociation propre au certificat apporte-t-elle après la clôture du marché d'origine~?
\end{enumerate}

\medskip\noindent\textbf{Partie III --- La simulation.}
\begin{enumerate}[resume]
\item Donner les erreurs pendant les heures d'ouverture et pendant les heures de fermeture, avec et sans le proxy.
\item Quelle est l'ampleur des écarts sur la plage de chevauchement, et à quelle fréquence dépassent-ils un aller-retour de 10~points de base~?
\item Quel est le seuil de conversion, et pourquoi est-il beaucoup plus élevé que l'écart courant~?
\item Comment la fourchette du teneur de marché doit-elle évoluer au fil de la séance~?
\end{enumerate}

\medskip\noindent\textbf{Partie IV --- Le verdict.}
\begin{enumerate}[resume]
\item Énoncer le \emph{résultat nommé}~: l'erreur de cotation du teneur de marché du certificat en heures fermées par rapport aux heures ouvertes, et le seuil de conversion après
frais.
\item Quand la conversion est-elle rentable~?
\item Quel effet les plafonds de détention étrangère ont-ils sur l'arbitrage~?
\item Laquelle des trois fiches de stratégie porte le risque de nuit~?
\item Comment couvrir la part propre à la société pendant la nuit~?
\item Pourquoi le backtest doit-il utiliser l'ouverture du marché d'origine, et non la clôture, comme sortie~?
\item Qu'apporterait un second proxy (un autre certificat issu du même marché)~?
\item En une phrase~: que vend le teneur de marché du certificat pendant les heures qui suivent la clôture du marché d'origine~?
\end{enumerate}
\end{problem}

\section{Questions d'entretien}

\begin{interviewq}[$\star$ \iqroles{trader}]\label{iq:hf:depositary-receipts-and-dual-listings:1}
Un certificat représente deux actions. Les actions cotent 800 pence et la livre sterling vaut 1.27~\$. Selon vous, où le certificat doit-il se négocier~?
\end{interviewq}

\begin{interviewq}[$\star\star$ \iqroles{chercheur}]\label{iq:hf:depositary-receipts-and-dual-listings:2}
Comment estimeriez-vous le fair value du certificat à 19:00, heure de Londres~?
\end{interviewq}

\begin{interviewq}[$\star\star$ \iqroles{trader}]\label{iq:hf:depositary-receipts-and-dual-listings:3}
Le certificat se négocie 1~\% au-dessus de la parité pendant une semaine. Pourquoi personne ne l'arbitre-t-il~?
\end{interviewq}

\begin{interviewq}[$\star\star$ \iqroles{risque}]\label{iq:hf:depositary-receipts-and-dual-listings:4}
Votre desk souhaite recevoir des certificats avant de livrer les actions au dépositaire. Que vérifiez-vous~?
\end{interviewq}

\begin{interviewq}[$\star\star$ \iqroles{développeur}]\label{iq:hf:depositary-receipts-and-dual-listings:5}
Concevez le service de fair value pour 300 certificats répartis sur quatre marchés d'origine et leurs devises.
\end{interviewq}

\begin{interviewq}[$\star\star\star$ \iqroles{chercheur}]\label{iq:hf:depositary-receipts-and-dual-listings:6}
Dérivez l'erreur quadratique moyenne du fair value en heures fermées en fonction du temps écoulé depuis la clôture du marché d'origine, avec et sans le proxy.
\end{interviewq}
